The flexor identification of record is the two-bracket fit of Section~\ref{sec:identification} on the four cross-validation tests: the $l_{0}=\qty{48.5}{\cm}$ and \qty{41.0}{\cm} tests in training and the $l_{0}=\qty{45.7}{\cm}$ and \qty{40.6}{\cm}-tendon tests held out. Its adopted terms are a length offset $\chi_{0} = \qty{+3.94}{\mm}$ and effective stiffnesses $\chi_{1} = \qty{39.98}{\kN\per\m}$ and $\chi_{2} = \qty{14.73}{\kN\per\m}$ (Table~\ref{tab:AdoptedXi}). At these terms the corrected model follows the measured torque across the range of motion, including the loss of torque near full extension that the rigid-body prediction does not produce (Figure~\ref{fig:FlxPin_optimized}). The fitting-set RMSE is \qty{2.165}{\newton\meter} for the $l_{0}=\qty{48.5}{\cm}$ test and \qty{1.084}{\newton\meter} for the $l_{0}=\qty{41.0}{\cm}$ test, and the held-out validation tests give \qty{1.368}{\newton\meter} at \qty{45.7}{\cm} and \qty{1.611}{\newton\meter} at \qty{40.6}{\cm} with tendon (Table~\ref{tab:FlxPinGof}). An initial single-bracket identification, fitted on the \qty{48.5}{\cm} test with the \qty{45.7}{\cm} test as validation, preceded the origin-side bracket and was superseded by this fit; its values are retained in Appendix~\ref{app:model_geometry}.
